% =====================================================================
% 5. MACHINE LEARNING MODELS
% =====================================================================
\section{Machine Learning models}

\subsection{Model selection}

The models were chosen to span the course from the most constrained to the most flexible, so that
the comparison answers a question (how much flexibility does this problem need?) rather than
searching for the highest score. For each family, the hyperparameters searched are those that control
its bias--variance trade-off \cite{james2021,hastie2009}.

\textbf{Linear and penalised linear models.} Multiple linear regression estimates
$y = \beta_0 + \sum_j \beta_j x_j + \varepsilon$ by least squares; with 0/1 predictors each
$\beta_j$ is the difference in mean age between patients with and without mutation $j$, holding the
others fixed. Its inference assumes independent, homoscedastic, normal errors, which are checked with
residual diagnostics. Ridge \cite{hoerl1970} and Lasso \cite{tibshirani1996} add the penalties
$\lambda\sum_j\beta_j^2$ and $\lambda\sum_j|\beta_j|$; Ridge shrinks all coefficients, while the
corners of the L1 constraint set some exactly to zero, i.e.\ Lasso selects variables. The penalty
$\lambda$ (\texttt{alpha}) is tuned, and inputs are standardised so that the penalty treats all
coefficients alike. Logistic regression models the log-odds of GBM as a linear function, with an L1
or L2 penalty and $C = 1/\lambda$ tuned. The penalty is also needed for a second reason: when a
predictor occurs in only one class, maximum-likelihood estimates do not exist (complete or
quasi-complete separation) \cite{albert1984}; NOTCH1 is such a predictor here.

\textbf{Naive Bayes} applies Bayes' rule assuming the predictors are independent given the class.
Age is discretised into quantile bins so that every column is Bernoulli; the number of bins and the
Laplace smoothing are tuned. The independence assumption is false for correlated mutations, which
often leaves the ranking of patients (and hence the AUC) reasonable while making the probabilities
over-confident \cite{hand2001}.

\textbf{Trees and ensembles.} A tree partitions the predictor space into boxes by recursive binary
splits chosen to reduce impurity (Gini index or entropy/deviance for classification, squared error for
regression) \cite{breiman1984}. A large tree overfits, so it must be pruned. The regression tree is
limited in advance by maximum depth and minimum leaf size; the classification tree is grown and then
pruned by cost-complexity pruning, which minimises impurity $+\ \alpha\cdot$(number of leaves) with
$\alpha$ (\texttt{ccp\_alpha}) tuned. Random forests average many trees grown on bootstrap samples,
trying a random subset of predictors (\texttt{max\_features}) at each split to decorrelate them
\cite{breiman2001}; averaging reduces variance, so the trees need no pruning, and the minimum leaf size
is tuned. Gradient boosting fits small trees sequentially to the gradient of the loss of the current
ensemble; XGBoost \cite{chen2016} is the implementation used, with the number of trees, the shrinkage
(learning rate) and the depth tuned.

\textbf{Support vector machines} find the maximum-margin hyperplane allowing violations penalised by
$C$, and the kernel trick gives non-linear boundaries; linear and RBF kernels are compared, with $C$
and the RBF reach $\gamma$ tuned \cite{cortes1995}. Probabilities for the AUC are obtained by Platt
scaling \cite{platt1999}. Support vector regression fits a function within an $\varepsilon$-tube,
ignoring errors smaller than $\varepsilon$ years.

\textbf{Baselines.} The mean age and the class proportions give the reference that any useful model
must beat.

\subsection{Complementary analyses}

\textbf{Multiclass classification.} Logistic regression extends to $K$ classes by giving every class
$k$ its own linear score $z_k = \beta_{k0} + \boldsymbol\beta_k^\top\mathbf x$ and turning the scores
into probabilities with the softmax function, $P(Y=k\mid\mathbf x) = e^{z_k}/\sum_{j} e^{z_j}$
(multinomial regression) \cite{james2021}. Trees and random forests handle several classes directly,
with the Gini index or the entropy computed over the $K$ class proportions of each node, and XGBoost
fits one tree per class at every boosting round on the softmax loss \cite{chen2016}. The main metric is
the multiclass AUC of Hand and Till, the mean over all pairs of classes $(i,j)$ of the AUC of
$P(Y=i)$ separating class $i$ from class $j$, averaged with the reverse comparison; like the binary AUC
it needs no threshold and does not depend on the class proportions \cite{handtill2001}. Balanced
accuracy (mean recall over the classes; 0.25 for a constant prediction with four classes) and the macro
F1 score describe the predicted class.

\textbf{SHAP values.} Permutation importance says how much performance is lost without a predictor,
but not in which direction a predictor pushes each prediction. SHAP values \cite{lundberg2017} split one
prediction $f(\mathbf x)$ into additive contributions, $f(\mathbf x) = \phi_0 + \sum_j \phi_j$, where
$\phi_0$ is the average prediction and $\phi_j$ is the Shapley value of predictor $j$
\cite{shapley1953}: the average change in the prediction when $j$ is added to every possible subset of
the other predictors. For trees and tree ensembles the values are computed exactly by TreeSHAP
\cite{lundberg2020}; for a linear model with independent predictors they reduce to
$\phi_j = \beta_j(x_j - \bar x_j)$. SHAP values describe the fitted model, not the disease: correlated
predictors share the credit in a way that depends on the algorithm, and a large value is not evidence of
a causal effect.

\textbf{Causal analysis.} A prediction model may use any variable that carries information; a causal
question requires deciding \emph{which} variables to adjust for, and that decision comes from a causal
diagram (a directed acyclic graph, DAG) rather than from the fit of the model
\cite{pearl2009,huntingtonklein2021}. On a path from an exposure to an outcome, a \emph{confounder} is a
common cause and must be adjusted for; a \emph{mediator} lies on the causal path, and adjusting for it
removes the part of the effect that goes through it; a \emph{collider} is a common effect, and adjusting
for it creates a spurious association. With an exposure $E$, a mediator $M$, an outcome $Y$ and the same
covariates, three least-squares regressions give the total effect $c$ ($Y$ on $E$), the effect $a$ of
$E$ on $M$, and the direct effect $c'$ and the effect $b$ of $M$ ($Y$ on $E$ and $M$); for ordinary
least squares on the same patients the total effect splits exactly as $c = c' + ab$, where $ab$ is the
indirect effect through the mediator \cite{baron1986,mackinnon1995}. The split is causal only under the
assumptions of the diagram, in particular that no unmeasured variable affects both the mediator and the
outcome.
